\documentclass[11pt,a4paper]{article}
\usepackage[margin=1in]{geometry}
\usepackage{amsmath}
\usepackage[numbers,sort&compress]{natbib}
\usepackage{hyperref}
\usepackage{booktabs}
\usepackage[swedish]{babel}

\title{W.B.\ Hutchinsons tekniska bidrag \\
       till den fysikaliska metallurgin}
\author{Teknisk sammanfattning}
\date{Mars 2026}

\begin{document}
\maketitle

\begin{abstract}
W.B.\ (Bevis) Hutchinson, som tillbringat st\"orre delen av sin
karri\"ar vid Svenska institutet f\"or metallforskning (SIMR, senare
Swerea KIMAB, numera Swerim) i Stockholm, har gjort grundl\"aggande
bidrag till f\"orst\aa{}elsen av kristallografisk textur,
rekristallisationsmekanismer och struktur--egenskapssamband i st\aa{}l
och aluminiumlegeringar.  Med \"over 180 publikationer och mer \"an
7\,000 citeringar sp\"anner hans arbete \"over kolfattiga pl\aa{}tst\aa{}l,
interstitiellt fria (IF) st\aa{}l, kornorienterade elst\aa{}l,
mikrolegerade konstruktionsst\aa{}l och aluminiumpl\aa{}tprodukter.
Detta dokument sammanfattar de viktigaste tekniska temana i hans
forskning.
\end{abstract}

%% ===================================================================
\section{Texturutveckling i kolfattiga pl\aa{}tst\aa{}l}
\label{sec:bcc-texture}

Hutchinsons \"oversiktsartikel fr\aa{}n 1984 i \emph{International
Materials Reviews}, \emph{``Development and control of annealing textures
in low-carbon steels''} \cite{Hutchinson1984IMR}, \"ar ett av hans mest
inflytelserika arbeten.  Den behandlar bildningen av
$\gamma$-fibertexturen ($\{111\}\langle uvw\rangle$) som utvecklas vid
rekristallisationsgl\"odgning av kallvalsad ferritisk st\aa{}lpl\aa{}t.

\subsection{$\gamma$-Fibern och dragbarhet}

Lankford-koefficienten $\bar{r}$ (normalt plastiskt anisotropif\"orh\aa{}llande)
styr djupdragbarheten hos st\aa{}lpl\aa{}t.  H\"oga $\bar{r}$-v\"arden
korrelerar med en stark $\gamma$-fiber, dvs.\ en \"overvikt av korn med
$\{111\}$-plan parallella med pl\aa{}tytan.  Den konkurrerande
texturkomponenten \"ar $\{100\}\langle 0vw\rangle$ ($\alpha$-fibern), som
f\"ors\"amrar dragbarheten.  Hutchinson visade att f\"orh\aa{}llandet
mellan $\{111\}$- och $\{100\}$-intensiteten i rekristallisations\-texturen
\"ar den avg\"orande faktorn f\"or $\bar{r}$-v\"ardet, och att detta
f\"orh\aa{}llande kan p\aa{}verkas genom legeringskemi (kol, kv\"ave,
mangan och mikrolegeringstills\"atser), varmbandsstruktur,
kallvalsningsreduktion och gl\"odgningspraxis (l\aa{}dgl\"odgning
\emph{vs.}\ kontinuerlig gl\"odgning).

\subsection{Nukleering fr\aa{}n $\gamma$-fiberkorn}

I deformerat kolfattigt st\aa{}l utvecklar
$\{111\}\langle 112\rangle$-korn t\"ata skjuvband under kallvalsning.
Hutchinson identifierade dessa skjuvband som preferentiella
nukleeringsst\"allen f\"or rekristalliserade korn som \"arver
$\{111\}$-orienteringen.  Denna ``orienterade nukleering'' fr\aa{}n
skjuvbandade $\gamma$-fiberkorn \"ar den huvudsakliga mekanismen genom
vilken en stark $\gamma$-fiberrekristallisationstextur uppr\"attas, vilket
leder till h\"oga $\bar{r}$-v\"arden.

\subsection{Interstitiellt fria (IF) st\aa{}l}

Utvecklingen av IF-st\aa{}l (stabiliserade med Ti och/eller Nb f\"or att
binda allt interstitiellt C och N) skapade ett idealiskt system f\"or
optimering av $\gamma$-fibertexturen.  Utan Cottrell-atmosf\"arer eller
fina karbidutskiljningar som st\"or kornr\"andernas r\"orlighet utvecklar
IF-st\aa{}l extremt starka $\{111\}$-rekristallisationstexter och
uppn\aa{}r $\bar{r}$-v\"arden \"over~2.  Hutchinsons texturanalys\-ramverk
ligger till grund f\"or den industriella processdesignen f\"or dessa
st\aa{}l.

%% ===================================================================
\section{Rekristallisationsmekanismer och substruktur}
\label{sec:rex}

Ett \aa{}terkommande tema i Hutchinsons arbete \"ar sambandet mellan
deformationssubstruktur och nukleering av rekristallisation.

\subsection{Deformationssubstrukturer}

I en \"oversikt publicerad i \emph{Materials Science Forum}
\cite{Hutchinson2007MSF} granskade Hutchinson de dislokationssubstrukturer
som skapas vid plastisk deformation av metaller, med s\"arskild
uppm\"arksamhet p\aa{} substrukturer som rymmer skarpa
misorientationer.  Dessa element med h\"og misorientation---deformationsband,
\"overg\aa{}ngsband, skjuvband och omr\aa{}den n\"ara kornr\"ander d\"ar
gitterrotationer ackumuleras---\"ar av grundl\"aggande betydelse
eftersom de tillhandah\aa{}ller de h\"ogvinkliga kornr\"ander som beh\"ovs
f\"or nukleering av rekristallisation.

Olika deformationsmoder och kristallorientringar ger upphov till olika
typer av substruktur, s\aa{} att nukleeringsmekanismen och den
resulterande rekristallisationstexturen beror p\aa{} den tidigare
deformationsv\"agen.  Denna insikt ger praktisk h\"avst\aa{}ngseffekt:
genom att kontrollera deformationsf\"orh\aa{}llandena (temperatur,
t\"ojning, t\"ojningsv\"ag) kan man p\aa{}verka vilka
nukleeringsmekanismer som \"ar aktiva och d\"armed styra den
rekristalliserade texturen.

\subsection{Debatten om orienterad nukleering kontra selektiv tillv\"axt}

Ursprunget till rekristallisationstexter har debatterats i decennier.
Tv\aa{} klassiska hypoteser \"ar: (i)~orienterad nukleering---nya korn
nukleerar f\"oretr\"adesvis fr\aa{}n vissa orienteringar; och (ii)~selektiv
tillv\"axt---korn med vissa orienteringar har en mobilitetsf\"ordel och
konsumerar den deformerade matrisen snabbare.  Hutchinsons experimentella
arbete, s\"arskilt med hj\"alp av HVEM (h\"ogsp\"anningselektronmikroskopi)
f\"or att observera nukleering in situ, bidrog med bevis f\"or orienterad
nukleering som den dominerande faktorn i m\aa{}nga system, \"aven om han
erk\"ande att b\aa{}da mekanismerna kan verka och att en kombinerad
modell \"ar n\"odv\"andig f\"or kvantitativ f\"oruts\"agelse.

%% ===================================================================
\section{Kubtextur i FCC-metaller}
\label{sec:cube}

I FCC-metaller (aluminium, koppar, nickel) producerar kraftig
kallvalsning en $\beta$-fibervalsningsextur, och efterf\"oljande
gl\"odgning utvecklar kubtexturen
$\{100\}\langle 001\rangle$ som den dominerande
rekristallisationskomponenten.

\subsection{Ridha--Hutchinson-mekanismen}

Ridha och Hutchinson \cite{Ridha1982} f\"oreslog en mekanism f\"or att
f\"orklara varf\"or kubinriktade regioner i den deformerade
mikrostrukturen fungerar som preferentiella nukleeringsst\"allen.  Vid
plant\"ojningsvalsning \"ar fyra glidsystem aktiva men med bara tv\aa{}
distinkta Burgersvektorer.  I kuborienterade band ger denna balanserade
glidaktivitet en relativt stabil substruktur med l\aa{}g
dislokationst\"athet.  Dessa kubband \aa{}terh\"amtar sig l\"att och
bildar subkorn begr\"ansade av h\"ogvinkliga kornr\"ander---f\"oruts\"att\-ningarna
f\"or livskraftiga rekristallisationskimer.  De resulterande kimerna
v\"axer snabbt in i den omgivande kraftigt deformerade
$\beta$-fibermatrisen, vilket f\"orklarar kubkomponentens dominans.

Denna modell har citerats brett och f\"orblir standardf\"orklaringen
f\"or kubtexturnukleering i valsade och gl\"odgade FCC-metaller.

%% ===================================================================
\section{Kornorienterade elst\aa{}l}
\label{sec:goss}

Hutchinson bidrog till f\"orst\aa{}elsen av sekund\"ar rekristallisation
och Goss-texturen $\{110\}\langle 001\rangle$ i Fe--Si-elst\aa{}l.  I
kornorienterat elst\aa{}l v\"axer en mycket liten andel
($\sim\!10^{-6}$) av de prim\"arrekristalliserade kornen med
Goss-orientering onormalt under en h\"ogtemperaturgl\"odgning och
konsumerar alla andra korn.  Denna ``sekund\"ara rekristallisation''
producerar en skarp Goss-textur som ger utm\"arkta magnetiska egenskaper
(l\aa{}ga k\"arnf\"orluster, h\"og permeabilitet l\"angs
valsningsriktningen).

Hutchinson och Homma \cite{Homma2003} unders\"okte orienteringsberoendet
hos sekund\"ar rekristallisation i kiselj\"arn och tillhandah\"oll
experimentella bevis f\"or vilka kornr\"andsmisorienteringar som
fr\"amjar onormal korntillv\"axt och d\"armed selekterar
Goss-orienteringen.  Han unders\"okte ocks\aa{} skjuvbandens bidrag
till ursprunget f\"or Goss-orienterade kimer i den prim\"ara
rekristallisationstexturen, och kopplade d\"armed sitt bredare arbete om
deformationssubstruktur till detta industriellt kritiska problem.

%% ===================================================================
\section{Mikrolegerade st\aa{}l och vanadin}
\label{sec:microalloy}

Vid SIMR medf\"orfattade Hutchinson en stor monografi om vanadinets roll
i mikrolegerade st\aa{}l \cite{Lagneborg2014}, tillsammans med
R.~Lagneborg, T.~Siwecki och S.~Zajac.

\subsection{Fysikalisk metallurgi f\"or vanadin i st\aa{}l}

Monografin j\"amf\"or systematiskt vanadin med andra
mikrolegeringselement (Nb, Ti) avseende termodynamisk stabilitet,
l\"oslighet i austenit och utskiljningsbeteende.  En avg\"orande
skillnad \"ar att V(C,N) har relativt h\"og l\"oslighet i austenit
j\"amf\"ort med NbC eller TiN.  Vid typiska slutvalsningstemperaturer
($>900\,{}^\circ$C) f\"orblir n\"astan allt vanadin i fast l\"osning,
medan Nb f\"aller ut rikligt i austenit vid dessa temperaturer.
F\"oljaktligen \"ar vanadinets huvudsakliga styrkebidrag genom
utskiljningsh\"ardning i ferrit, snarare \"an genom
austenitkornf\"orfining (som \"ar Nb:s och Ti:s dom\"an).

\subsection{Utskiljningsmekanismer i ferrit}

Monografin identifierar tre distinkta utskiljningstyper f\"or V(C,N) i
ferrit:

\begin{enumerate}
  \item \textbf{Interfasutskiljning} sker vid den migrerande
    $\gamma/\alpha$-gr\"ansytan under austenit-till-ferrit-omvandlingen.
    Rader av utskiljningar bildas vid successiva positioner av
    gr\"ansytan, vilket producerar de karakt\"aristiska plan\"ara
    m\"onstren som observeras i TEM.  Lagneborg och Zajac
    \cite{LagneborgZajac2001} f\"oreslog en stegmigrationsmodell d\"ar
    V(C,N) nukleerar p\aa{} station\"ara terrasser p\aa{} gr\"ansytan;
    n\"ar stegen sveper \"over l\"amnas den utskiljningsbelastade
    terrassen kvar och en ny terrass exponeras f\"or n\"asta generation
    av nukleering.  Avståndet mellan skikten styrs av stegh\"ojden och
    omvandlingstemperaturen.

  \item \textbf{Allm\"an (slumpm\"assig) utskiljning} sker efter att
    omvandlingen \"ar fullbordad, under l\aa{}ngsam avkylning eller
    isoterm h\aa{}llning i ferriten.  Till skillnad fr\aa{}n
    interfasutskiljning nukleerar partiklarna p\aa{} dislokationer och
    subkornr\"ander, vilket ger en slumpm\"assig rumslig f\"ordelning.

  \item \textbf{Fibr\"os utskiljning} kan intr\"affa n\"ar gr\"ansytan
    r\"or sig mycket l\aa{}ngsamt, varvid utl\"angda
    utskiljningsfibrer bildas snarare \"an diskreta partiklar.
\end{enumerate}

Fr\aa{}n Orowan-sambandet h\"arleder monografin kvantitativa
styrkeskattningar: approximately 6\,MPa per 0,001\,vikt\%\,N och
5,5\,MPa per 0,01\,vikt\%\,C f\"or interfasutskild V(C,N).  Ett
viktigt resultat \"ar att kol och kv\"ave bidrar synergistiskt till
utskiljningsh\"ardningen.  Vid h\"ogre kv\"avehalter \"ar
V(C,N)-partiklarna mindre och t\"atare f\"ordelade eftersom VN har en
l\"agre l\"oslighetsprodukt \"an VC, vilket leder till en h\"ogre
nukleeringshastighet.  Den finare dispersionen inkorporerar d\aa{} ocks\aa{}
mer kol, s\aa{} att kombinerade V--N--C-st\aa{}l uppn\aa{}r h\"ogre
utskiljningsh\"ardning \"an summan av V--N- och V--C-bidragen var f\"or
sig.

\subsection{VN som ferritnukleator}

Monografin beskriver en mekanism f\"or ferritkornf\"orfining genom
intragranul\"ar nukleering p\aa{} VN-utskiljningar i austenit.
VN-partiklar som bildas under avkylning genom austenitomr\aa{}det (eller
under ett \"overf\"orings-/\aa{}teruppv\"armningssteg) antar
Baker--Nutting-orienteringsrelationen med ferrit,
$(100)_\text{VN}\parallel(100)_\alpha$,
$[011]_\text{VN}\parallel[001]_\alpha$.  Denna l\aa{}generigr\"ansyta
g\"or VN till en effektiv heterogen nukleeringsplats f\"or ferrit, vilket
fr\"amjar intragranul\"ar ferritbildning och ger ett
omvandlingsf\"orh\aa{}llande $D_\gamma/D_\alpha$ s\aa{} h\"ogt som 5--10
(j\"amf\"ort med $\sim$2 f\"or konventionell kornr\"andsnukleering).
TiN-partiklar kan fungera som f\"orl\"opare f\"or VN-utskiljning via
samutskiljning, vilket f\"orst\"arker denna mekanism i Ti--V--N-st\aa{}l.

\subsection{Rekristallisationskontrollerad valsning (RCR)}

Den h\"oga l\"osligheten av V(C,N) i austenit inneb\"ar att vanadin inte
f\"ordr\"ojer rekristallisation mellan valspass i samma utstr\"ackning
som niob.  F\"oljaktligen l\"ampar sig vanadinmikrolegering inte f\"or
konventionell kontrollerad valsning (CR) under
rekristallisationsstopptemperaturen.  I st\"allet utvecklade Hutchinson
och medarbetare konceptet \emph{rekristallisationskontrollerad valsning}
(RCR) \cite{Zajac1991,Siwecki1995}:

\begin{itemize}
  \item Valsningen utf\"ors vid relativt h\"oga slutvalsningstemperaturer
    ($\geq 950\,{}^\circ$C), d\"ar austeniten rekristalliserar fullst\"andigt
    mellan passen och uppr\"atth\aa{}ller en fin likaxlig kornstruktur
    genom upprepade rekristallisationscykler.
  \item Sm\aa{} tills\"atser av Ti ($\sim$0,01\,\%) tillhandah\aa{}ller
    TiN-partiklar som f\"astpinnar austenitkornr\"anderna via
    Zener-motst\aa{}nd och f\"orhindrar kornf\"orgrovning efter
    rekristallisation.
  \item Accelererad avkylning (ACC) efter valsning, f\"oljd av
    upphaspling vid 550--$600\,{}^\circ$C f\"or band, driver
    $\gamma\to\alpha$-omvandlingen till fina ferrit--bainit-strukturer
    och bevarar vanadin i l\"osning f\"or fin V(C,N)-utskiljning under
    den l\aa{}ngsamma avkylningen i rullen.
\end{itemize}

Datormodellen MICDEL, utvecklad vid Svenska institutet f\"or
metallforskning, anv\"andes f\"or att optimera RCR-passscheman genom att
simulera den kopplade utvecklingen av austenitkornstorlek,
rekristallisationsandel och utskiljningsf\"ordelning under varmvalsning
\cite{Siwecki1992}.  Ett kritiskt praktiskt resultat var att sm\aa{}
``kalibreringspass'' i slutet av valssekvensen m\aa{}ste undvikas: l\aa{}ga
t\"ojningar vid h\"og temperatur kan orsaka onormal kornf\"orgrovning
snarare \"an f\"orfining, eftersom rekristallisationsnukleeringen \"ar
otillr\"ackligt riklig (fenomenet ``kritisk t\"ojningsgl\"odgning'').

\subsection{Kv\"avets inverkan}

Ett centralt tema i monografin \"ar kv\"avets kraftfulla och ofta
underskattade roll i V-mikrolegerade st\aa{}l.  \"Okning av kv\"avehalten
fr\aa{}n $\sim$0,003\,\% till $\sim$0,013\,\% i ett st\aa{}l med
0,09\,\%V h\"ojer str\"ackgr\"ansen med $\sim$120\,MPa f\"or
RCR-bearbetat material \cite{Zajac1991}.  Detta h\"arr\"or fr\aa{}n tv\aa{}
\"omsesidigt f\"orst\"arkande mekanismer: (i)~t\"atare VN-nukleering
ger finare utskiljningsdispersioner med mindre partikelavst\aa{}nd, vilket
f\"orb\"attrar Orowan-h\"ardningen; (ii)~VN-partiklar i austenit
fungerar som intragranul\"ara ferritnukleatorer
(avsnitt~\ref{sec:microalloy}), vilket f\"orfinar ferritkornstrukturen.
Monografin visar att kv\"ave \"ar ett v\"asentligt mikrolegeringselement
tillsammans med vanadin---inte bara en skadlig f\"ororening som man
traditionellt ansett.

%% ===================================================================
\section{Aluminiumpl\aa{}t: textur, \"oronbildning och formbarhet}
\label{sec:aluminium}

Hutchinson till\"ampade sin texturexpertis p\aa{} bearbetning av
aluminiumpl\aa{}t.  Arbeten med Oscarsson och Ekstr\"om
\cite{Oscarsson1991} unders\"okte inverkan av utg\aa{}ngsmikrostrukturen
p\aa{} textur och \"oronbildning i aluminiumpl\aa{}t efter kallvalsning
och gl\"odgning.

\subsection{Kontroll av \"oronbildning}

``\"Oronbildning''---bildningen av en v\aa{}gig kant (\"oron) vid
djupdragning av en cirkul\"ar pl\aa{}trondell---\"ar en direkt
manifestation av plan anisotropi orsakad av kristallografisk textur.
En valsningstextur ger $0^\circ/90^\circ$-\"oron; en
kubrekristallisationstextur ger $45^\circ$-\"oron.  Genom att balansera
de relativa styrkorna hos dessa tv\aa{} texturkomponenter via kontroll
av slutvalsningstemperaturen vid varmvalsning, graden av
rekristallisation f\"ore kallvalsning, kallvalsningsreduktion och
gl\"odgningsf\"orh\aa{}llanden kan \"oronbildningen minimeras.

\subsection{Texturmodellering genom hela processen}

Hutchinsons f\"orst\aa{}else av texturutvecklingsmekanismerna---bildning
av deformationstextur, \aa{}terh\"amtning, nukleering och
korntillv\"axt---gav den fysikaliska grunden f\"or
genomg\aa{}endetexturmodeller till\"ampade p\aa{} produktion av
aluminiumpl\aa{}t, och kopplade processparametrar till sluttextur och
d\"armed till mekanisk anisotropi, \"oronbeteende och formbarhet.

%% ===================================================================
\section{Laserultraljud f\"or processtyring i realtid}
\label{sec:lus}

I senare arbeten unders\"okte Hutchinson laserultraljudstekniker (LUS)
f\"or ber\"oringsfri realtidsm\"atning av mikrostruktur under
varmvalsning av st\aa{}l \cite{Hutchinson2011LUS}.  F\"ors\"ok vid
SSAB:s varmbandsvalsverk i Borl\"ange visade att lasergenererade och
laserdetekterade ultraljudsv\aa{}gor kan m\"ata
austenit-till-ferrit-omvandlingstillst\aa{}ndet, rekristallisationsandelen
och kornstorleken p\aa{} ett r\"orligt band vid utloppet.  Detta arbete
representerar en frontlinje inom sluten mikrostrukturstyrning f\"or
industriell st\aa{}lproduktion, och kopplar Hutchinsons djupa f\"orst\aa{}else
av fasomvandlingar och rekristallisationskinetik till praktisk
process\"overvakning.

%% ===================================================================
\section{Deformationsmikrostrukturer och omvandlingstexter}
\label{sec:deformation}

En artikel fr\aa{}n 1999 i \emph{Philosophical Transactions of the Royal
Society} \cite{Hutchinson1999Phil} gav en omfattande behandling av
deformationsmikrostrukturer och texter i st\aa{}l.  Detta arbete
f\"orenade Hutchinsons observationer om hur olika
deformationsmoder---plant\"ojning (valsning), enaxlig dragning,
skjuvning---skapar karakt\"aristiska dislokationssubstrukturer i
BCC-j\"arn, och hur dessa substrukturer i sin tur styr de texter som
utvecklas vid efterf\"oljande rekristallisation eller fasomvandling.
Tillsammans med medforfattarna Ryde och Bate studerade han ocks\aa{}
omvandlingstexter fr\aa{}n austenit-till-ferrit-omvandlingen och
unders\"okte hur Kurdjumov--Sachs- och
Nishiyama--Wassermann-orienteringsrelationerna mellan moderausteint och
produktferrit p\aa{}verkar texturen i den slutliga ferritiska
mikrostrukturen.

%% ===================================================================
\section{Anisotropi i metaller: en syntes}
\label{sec:anisotropy}

En kritisk bed\"omning fr\aa{}n 2015 i \emph{Materials Science and
Technology}, \emph{``Anisotropy in metals''} \cite{Hutchinson2015},
fungerar som en avslutande syntes av Hutchinsons karri\"arteman.  Den
behandlar elastisk och plastisk anisotropi i b\aa{}de BCC- och
FCC-system och visar hur kristallografisk textur styr riktningsberoende
mekaniska egenskaper.  Bed\"omningen omfattar enkristallens elastiska
konstanter, polykristallmedelvärdesmodeller (Voigt, Reuss, Hill,
Kr\"oner), sambandet mellan ODF (orienteringsf\"ordelningsfunktionen) och
$r$-v\"arden, samt praktiska implikationer f\"or pl\aa{}tformning.  Den
demonstrerar den enande tr\aa{}den i Hutchinsons arbete: textur \"ar bron
mellan bearbetning, mikrostruktur och mekanisk prestanda.

I en artikel fr\aa{}n 2025 \cite{Hutchinson2025pssb} fortsatte
Hutchinson detta tema genom att experimentellt m\"ata longitudinella
elastiska v\aa{}ghastigheter i polykristallint j\"arn fr\aa{}n
rumstemperatur till $880\,{}^\circ\mathrm{C}$ och j\"amf\"ora
resultaten med Hill- och Kr\"oner-medelvärdesmodellerna.

%% ===================================================================
\section{\"Ovriga anm\"arkningsv\"arda arbeten}
\label{sec:other}

\subsection{L\"arobok: An Introduction to Textures in Metals}

Hutchinson medforfattade med M.~Hatherly monografin \emph{An
Introduction to Textures in Metals} (Institute of Metals, London, 1979;
2:a~uppl.\ 1990) \cite{Hatherly1979}.  Denna l\"arobok har fungerat som
en standardreferens f\"or omr\aa{}det kristallografisk textur i metaller.

\subsection{Gibeon-meteoriten}

I en uppfinningsrik till\"ampning av EBSD-teknik
(elektron\aa{}terspriddningsdiffraktion) studerade Hutchinson och
Hagstr\"om \cite{Hutchinson2006Met} Widmanst\"atten-ferritomvandlingen i
Gibeon-j\"arnmeteoriten och fann orienteringsrelationer mellan
ferritplattorna och moderausteniten f\"orenliga med b\aa{}de
Nishiyama--Wassermann- och Kurdjumov--Sachs-varianter---ett naturligt
fasomvandlingsexperiment som \"agt rum \"over geologiska tidsskalor.

\subsection{Partiell gl\"odgning av kallvalsade st\aa{}l}

Tillsammans med R.K.~Ray och C.~Ghosh \cite{Ray2011} granskade
Hutchinson partiell gl\"odgning (\emph{back-annealing})---v\"armebehandling
av kallvalsade st\aa{}l avsedd att uppn\aa{} partiell \aa{}terh\"amtning
och/eller partiell rekristallisation snarare \"an fullst\"andig
rekristallisation.  Denna processv\"ag m\"ojligg\"or skr\"addarsydda
avv\"agningar mellan styrka och duktilitet genom att kontrollera
andelen \aa{}terh\"amtad kontra rekristalliserad mikrostruktur.

%% ===================================================================
\section{Vanadinmonografin: en utf\"orlig sammanfattning}
\label{sec:vanadium-monograph}

Det mest omfattande enskilda verket d\"ar Hutchinsons metallurgiska
t\"ankande \"ar framlagt \"ar den 101 sidor l\aa{}nga monografin
\emph{The Role of Vanadium in Microalloyed Steels} \cite{Lagneborg2014},
f\"orfattad med R.~Lagneborg, T.~Siwecki och S.~Zajac vid Svenska
institutet f\"or metallforskning (numera Swerim).  Ursprungligen
publicerad 1999 i \emph{Scandinavian Journal of Metallurgy} utkom en
ut\"okad upplaga som Swerea KIMAB-rapport KIMAB-2014-115, sponsrad av
Vanadium International Technical Committee (Vanitec).  Monografin \"ar
tillagnad minnet av Michael Korchynsky, en pioneer inom
vanadinmikrolegering.  Det som f\"oljer \"ar en frist\aa{}ende
sammanfattning av detta verk, skriven f\"or en vetenskapligt utbildad
l\"asare som kanske inte \"ar specialist inom metallurgi.

\subsection{Bakgrund: vad \"ar ett mikrolegerat st\aa{}l?}

Konstruktionsst\aa{}l---anv\"anda i byggnader, broar, fartyg,
r\"orledningar och fordon---m\aa{}ste kombinera h\"og styrka med god
duktilitet och seghet (motst\aa{}nd mot pl\"otsligt, spr\"ott brott).
Det traditionella s\"attet att f\"orst\"arka st\aa{}l \"ar att
till\"agga mer kol, men kol g\"or st\aa{}let spr\"ott och sv\aa{}rt att
svetsa.  \emph{Mikrolegerade} st\aa{}l (eller HSLA,
\emph{High-Strength Low-Alloy}) l\"oser detta dilemma genom att
till\"agga mycket sm\aa{} m\"angder---typiskt 0,01--0,15\,\%---av
\"amnen s\aa{}som vanadin (V), niob (Nb) och titan (Ti).  Dessa
``mikrotills\"atser'' uppn\aa{}r sin effekt inte genom
fastl\"osningsh\"ardning (som kr\"aver stora m\"angder) utan genom tv\aa{}
mer effektiva mekanismer:
\begin{enumerate}
  \item \textbf{Kornf\"orfining:} att g\"ora de kristallina kornen i
    st\aa{}let finare, vilket samtidigt \"okar b\aa{}de styrkan och
    segheten---en s\"allsynt kombination i materialvetenskap, eftersom
    dessa egenskaper vanligen avv\"ager mot varandra.
  \item \textbf{Utskiljningsh\"ardning:} att bilda enorma antal
    nanometerskaliga karbid- och nitridpartiklar som hindrar
    dislokationernas r\"orelse (de linjedefekter som b\"ar plastisk
    deformation), vilket h\"ojer str\"ackgr\"ansen.
\end{enumerate}

De tre mikrolegeringselementen (V, Nb, Ti) skiljer sig fundamentalt i
\emph{n\"ar} och \emph{var} deras partiklar bildas, och monografin \"ar
strukturerad kring denna distinktion.

\subsection{Den termodynamiska grunden: l\"osligheten avg\"or}

St\aa{}l bearbetas genom uppv\"armning till h\"og temperatur d\"ar det
existerar i den kubiskt ytcentrerade (FCC) kristallstrukturen kallad
\emph{austenit} ($\gamma$), varmvalsas sedan, och slutligen kyls s\aa{}
att det omvandlas till den kubiskt rymdcentrerade (BCC) fasen kallad
\emph{ferrit} ($\alpha$) som utg\"or slutprodukten.  Den avg\"orande
fr\aa{}gan f\"or varje mikrolegeringselement \"ar: hur mycket av det
l\"oses i austeniten vid valstemperaturen, och hur mycket har redan
fallit ut?

\begin{itemize}
  \item \textbf{TiN} \"ar extremt stabilt och n\"astan ol\"osligt i
    austenit.  Det f\"aller ut under gjutning och best\aa{}r genom alla
    efterf\"oljande bearbetningssteg.  Dess roll \"ar att f\"astpinna
    austenitkornr\"ander (Zener-pinning), vilket f\"orhindrar
    kornf\"orgrovning vid \aa{}teruppv\"armning och svetsning.
  \item \textbf{NbC} har mellanliggande stabilitet.  Det f\"aller ut i
    austenit under varmvalsning, s\"arskilt under
    $\sim$950\,${}^\circ$C.  Nb i fast l\"osning och som utskiljningar
    f\"ordr\"ojer b\aa{}de kraftigt austenitens rekristallisation, vilket
    ger utdragna (``pannkaks-'') austenitkorn som omvandlas till mycket
    fin ferrit.  Dock f\"aller det mesta Nb ut i austeniten, s\aa{} att
    relativt lite \aa{}terst\aa{}r f\"or utskiljningsh\"ardning i ferrit.
  \item \textbf{V(C,N)} har h\"og l\"oslighet i austenit.  Vid typiska
    slutvalsningstemperaturer f\"orblir n\"astan allt vanadin l\"ost.
    Det f\"aller ut f\"orst under eller efter
    $\gamma\to\alpha$-omvandlingen och bildar nanometerskaliga
    V(C,N)-partiklar i ferrit som ger stark utskiljningsh\"ardning.
\end{itemize}

Dessa skillnader \"ar inte brister---de \"ar komplement\"ara.
Monografin argumenterar f\"or att V, Nb och Ti b\"or ses som att de
fyller olika roller i en samordnad legeringsdesign.

\subsection{Tre utskiljningstyper i ferrit}

N\"ar austenit omvandlas till ferrit vid avkylning blir det vanadin som
var l\"ost i austeniten \"overm\"attat i ferriten och f\"aller ut.
Monografin identifierar tre typer:

\begin{enumerate}
  \item \textbf{Interfasutskiljning.}  N\"ar
    $\gamma/\alpha$-omvandlingsfronten sveper genom st\aa{}let
    nukleerar rader av V(C,N)-partiklar p\aa{} sj\"alva gr\"ansytan.
    Gr\"ansytan avancerar genom en stegmekanism: plana terrasser
    ackumulerar utskiljningar, medan laterala steg sveper \"over och
    l\"amnar den utskiljningsbelastade terrassen kvar och exponerar en
    ny yta f\"or n\"asta rad.  I transmissionselektronmikroskop ger
    detta sl\aa{}ende m\"onster av parallella utskiljningsskikt,
    med avst\aa{}nd p\aa{} 10--30\,nm.

  \item \textbf{Allm\"an (slumpm\"assig) utskiljning.}  Efter att
    omvandlingsfronten har passerat driver kvarvarande \"overm\"attnad
    ytterligare utskiljning p\aa{} dislokationer och subkornr\"ander i
    ferriten.  Dessa partiklar \"ar slumpm\"assigt f\"ordelade snarare
    \"an ordnade i skikt.

  \item \textbf{Fibr\"os utskiljning.}  Vid mycket l\aa{}ga
    gr\"ansytshastigheter kan utskiljningar v\"axa som kontinuerliga
    fibrer snarare \"an diskreta partiklar.  Detta \"ar mindre vanligt
    under industriella f\"orh\aa{}llanden.
\end{enumerate}

B\aa{}de interfas- och allm\"an utskiljning bidrar till
Orowan-h\"ardningsbidraget (\"okningen i str\"ackgr\"ans orsakad av
partiklar som dislokationer m\aa{}ste b\"oja sig runt).  Monografin ger
kvantitativa skattningar: approximately 6\,MPa per 0,001\,vikt\%\,N
och 5,5\,MPa per 0,01\,vikt\%\,C.

\subsection{Kv\"ave: fr\aa{}n f\"ororening till legeringselement}

Ett av monografins mest sl\aa{}ende argument \"ar rehabiliteringen av
kv\"ave som ett \emph{f\"ordelaktigt} legeringselement.  Konventionellt
anses kv\"ave i st\aa{}l vara skadligt: ``fritt'' kv\"ave orsakar
t\"ojnings\aa{}ldring och f\"orspr\"odning.  Men i vanadininneh\aa{}llande
st\aa{}l \"ar kv\"avet bundet som VN, vilket \"ar termodynamiskt mer
stabilt \"an VC.  Den h\"ogre nukleeringshastigheten hos VN ger finare,
t\"atare utskiljningsdispersioner, vilket dramatiskt f\"orb\"attrar
h\"ardningen.  Monografin visar att en f\"ordubbling av kv\"avehalten i
ett armeringsst\aa{}l med 0,11\,\%V (fr\aa{}n $\sim$85\,ppm till
$\sim$180\,ppm) h\"ojer str\"ackgr\"ansen med n\"astan
120\,MPa---en anm\"arkningsv\"ard effekt fr\aa{}n ett sp\aa{}r\"amne.

Dessutom verkar kol och kv\"ave synergistiskt.  De fina VN-partiklar som
nukleerar f\"orst tj\"anar som mallar f\"or efterf\"oljande V(C,N)-tillv\"axt
som inkorporerar kol.  Den kombinerade h\"ardningen \"overtr\"affar
vad vardera elementet skulle bidra ensamt.

\subsection{Ferritnukleering p\aa{} VN: att anv\"anda austenitomvandlingen
  som ett f\"orfiningverktyg}

Ett s\"arskilt elegant bidrag \"ar konceptet att anv\"anda
VN-utskiljningar som intragranul\"ara ferritnukleatorer.  Normalt
nukleerar ferritkorn vid austenitkornr\"ander under avkylning; den
resulterande ferritkornstorleken begr\"ansas av den tidigare
austenitkornstorleken.  Men VN-partiklar som f\"aller ut inuti
austenitkorn (under en avkylnings-/\aa{}teruppv\"armningscykel eller
l\"angvarig h\aa{}llning) har en kristallografisk orienteringsrelation
med ferrit (Baker--Nutting-relationen) som g\"or dem till kraftfulla
heterogena nukleeringsplatser.  Ferritkorn nukleerar inuti
austenitkornen p\aa{} dessa partiklar, vilket ger ett
omvandlingsf\"orfiningsf\"orh\aa{}llande $D_\gamma/D_\alpha$ p\aa{}
5--10 (j\"amf\"ort med $\sim$2 f\"or enbart konventionell
kornr\"andsnukleering).  Industriella f\"ors\"ok vid Riva Acciaio Sellero
(Italien) p\aa{} tunga balkar bekr\"aftade kornstorlek\-sreduktioner
fr\aa{}n 15\,$\mu$m till 10\,$\mu$m och str\"ackgr\"ans\"okningar p\aa{}
$\sim$130\,MPa.

\subsection{Bearbetningsv\"agar: RCR, CR och kylningens roll}

Monografin \"agnar betydande uppm\"arksamhet \aa{}t hur vanadinst\aa{}l
b\"or varmvalsas.  Det finns tv\aa{} huvudv\"agar:

\paragraph{Kontrollerad valsning (CR).} Den konventionella metoden:
slutvalsning vid l\aa{}ga temperaturer (under $\sim$900\,${}^\circ$C) f\"or
att deformera austeniten utan rekristallisation, varvid utdragna
``pannkaks''-korn skapas vars stora kornr\"andsarea ger fin ferrit.
Detta kr\"aver h\"oga valsningskrafter och l\"ampar sig f\"or
Nb-st\aa{}l, som starkt motst\aa{}r rekristallisation.

\paragraph{Rekristallisationskontrollerad valsning (RCR).} Den metod
som f\"oredras f\"or vanadinst\aa{}l: slutvalsning vid h\"ogre
temperaturer ($>$\,950\,${}^\circ$C), d\"ar austeniten rekristalliserar
fullst\"andigt mellan passen.  Kornf\"orfining uppn\aa{}s genom upprepad
rekristallisation snarare \"an pannkaksbildning.  TiN-partiklar
f\"orhindrar kornf\"orgrovning mellan passen.  RCR kr\"aver l\"agre
valsningskrafter (cirka 25\,\% mindre \"an CR f\"or samma produkt), ger
h\"ogre valsverksproduktivitet och \"ar den naturliga v\"agen f\"or
vanadin eftersom den h\"oga V(C,N)-l\"osligheten inneb\"ar att det inte
finns n\aa{}gra utskiljningar som f\"astpinnar austeniten mot
rekristallisation.

Efter valsning omvandlar \emph{accelererad avkylning} p\aa{}
utloppsbordet den fina austeniten till en \"annu finare
ferrit--bainit-blandning.  I bandprodukter hasplas bandet sedan vid
550--600\,${}^\circ$C, d\"ar l\aa{}ngsam avkylning i rullen ger ideala
f\"orh\aa{}llanden f\"or V(C,N)-utskiljning.  Monografin visar att
denna RCR+ACC-v\"ag uppn\aa{}r styrke- och seghetsniv\aa{}er j\"amf\"orbara
med eller b\"attre \"an konventionell CR, samtidigt som den \"ar
industriellt mer effektiv.

\subsection{Industriella till\"ampningar}

Kapitel~7 och~8 i monografin \"overblickar specifika
produkttill\"ampningar:

\paragraph{Grovpl\aa{}t.}  Fullskaliga f\"ors\"ok p\aa{}
0,01\,\%Ti--0,08\,\%V-pl\aa{}tar vid Pohang (Korea), Rautaruukki
(Finland) och Baoshan (Kina) uppn\aa{}dde str\"ackgr\"anser \"over
490\,MPa med utm\"arkt slagsseghet
(ITT$_{40\text{J}} < -80\,{}^\circ$C) via RCR+ACC-bearbetning.

\paragraph{Varmvalsad bandpl\aa{}t.}  Bandvalsning \"ar n\"astan idealt
l\"ampad f\"or V--N-mikrolegering: snabb vattenkylning p\aa{}
utloppsbordet bevarar vanadin i l\"osning, medan l\aa{}ngsam
rullkylning ger gott om tid f\"or fin V(C,N)-utskiljning.
Str\"ackgr\"anser p\aa{} 550\,MPa har uppn\aa{}tts i 8--10\,mm band med
slagomvandlingstemperaturer under $-40\,{}^\circ$C.

\paragraph{Bainitiskt bandst\aa{}l.}  En anm\"arkningsv\"ard utveckling
vid Swerea KIMAB och SSAB var ett V-mikrolegerat bainitiskt st\aa{}l med
m\aa{}let 700\,MPa str\"ackgr\"ans f\"or lastbilschassitill\"ampningar
\cite{Hutchinson2014Bainite}.  I detta system p\aa{}verkar vanadin inte
sj\"alva bainitomvandlingen utan h\"ammar \aa{}terh\"amtning av
dislokationsstrukturen under hasplingen, vilket bevarar styrkan \"aven
vid hasplingstemperaturer under 500\,${}^\circ$C---en praktisk f\"ordel
med tanke p\aa{} sv\aa{}righeten att exakt styra hasplingstemperaturen
p\aa{} industriella utloppsbord.

\paragraph{Armering (armeringsj\"arn).}  Armering \"ar den enskilt
st\"orsta till\"ampningen av V-mikrolegering r\"aknat i tonnage.  Med
h\"ogre kolhalter ($\sim$0,16--0,22\,\%) och h\"oga
slutvalsningstemperaturer f\"orblir allt vanadin i l\"osning f\"or
utskiljningsh\"ardning.  Monografin visar att 5~delar N per vikt
motsvarar 1~del V i termer av h\"ardning.  V--N-mikrolegering har blivit
den f\"oredragna v\"agen f\"or h\"ogh\aa{}llfasta armeringsj\"arn utanf\"or
Europa, med produktion s\"arskilt stor i Kina ($\sim$130\,Mton/\aa{}r
armering 2010).

\paragraph{Smidesdetaljer.}  V-mikrolegerade smidesst\aa{}l med
ferrit--perlitstruktur (t.ex.\ 30MNV56) erbjuder ett konkurrenskraftigt
alternativ till segh\"ardade (Q-T) sorter f\"or fordonskomponenter som
vevaxlar och vevstakar.  De eliminerar v\"armebehandlingssteg
(h\"ardning, anl\"opning, riktning), minskar bearbetningskostnaderna med
28\,\% och uppvisar j\"amf\"orbar utmattningsh\aa{}llfasthet, om \"an med
n\aa{}got l\"agre seghet.

\paragraph{S\"oml\"osa r\"or.}  En innovativ till\"ampning utnyttjar
intragranul\"ar VN-ferritnukleering f\"or att eliminera det
normaliseringsgl\"odgningssteg i linjen som traditionellt kr\"avs f\"or
att f\"orfina den grova mikrostrukturen som produceras vid
h\"ogtemperaturh\aa{}ltagning och valsning.

\subsection{Svetsbarhet}

Kapitel~8 tar upp en avg\"orande fr\aa{}ga: komprometterar den h\"ogre
kv\"avehalt som \"ar f\"ordelaktig f\"or styrkan svetsbarheten?  Den
v\"armep\aa{}verkade zonen (VPZ) intill en svets \aa{}teruppv\"arms till
n\"ara sm\"altpunkten, varvid utskiljningar l\"oses upp och
austenitkornstrukturen f\"orgrovnas.  Monografin visar att:

\begin{itemize}
  \item Sm\aa{} Ti-tills\"atser ($\sim$0,01\,\%) \"ar n\"odv\"andiga:
    TiN-partiklar \"overlever svetstemperaturer och begr\"ansar
    austenitkornf\"orgrovning i VPZ, vilket bevarar segheten vid
    smaltlinjen.
  \item F\"or svetsning med l\aa{}g till medeh\"og v\"armetillf\"orsel
    ($\Delta t_{8/5} < 30$\,s) uppvisar st\aa{}l med h\"ogre kv\"avehalt
    faktiskt \emph{b\"attre} grovkornig VPZ-seghet \"an varianter med
    l\aa{}gt kv\"ave, eftersom V och N i l\"osning fr\"amjar acikul\"ar
    ferrit snarare \"an grov kornr\"andsferrit, vilket ger en finare
    effektiv kornstorlek.
  \item F\"or h\"og v\"armetillf\"orsel ($>$10\,kJ/mm) kan segheten
    f\"ors\"amras i st\aa{}l med h\"ogt N p\aa{} grund av grov
    polygonal kornr\"andsferrit; detta mildras genom
    Ti/REM/Ca-mikrotills\"atser.
  \item Flerpassvetsning \emph{f\"orb\"attrar} faktiskt VPZ-segheten:
    VN som l\"osts upp av det f\"orsta passet \aa{}terutf\"alls under
    den termiska cykeln fr\aa{}n efterf\"oljande pass, varvid
    mikrostrukturen f\"orfinas.
\end{itemize}

Den \"overgripande slutsatsen \"ar att V--N-mikrolegerade st\aa{}l \"ar
fullt f\"orenliga med god svetsbarhet n\"ar vettiga gr\"anser f\"or
v\"armetillf\"orsel iakttas och sm\aa{} Ti-tills\"atser inkluderas.

\subsection{Monografin i sitt sammanhang}

Monografin av Lagneborg--Hutchinson--Siwecki--Zajac intar en unik
position i litteraturen om mikrolegerade st\aa{}l.  De n\"armaste
motsvarigheterna \"ar Gladmans \emph{The Physical Metallurgy of
Microalloyed Steels} (1997), som ger en bredare men mindre detaljerad
\"oversikt av alla mikrolegeringselement, och konferensrapporterna fr\aa{}n
Microalloying-serien (1975, 1995, 2007) som samlar kortare bidrag
fr\aa{}n m\aa{}nga forskargrupper.  Vanadinmonografin uts\"aker sig genom:
\begin{itemize}
  \item Sitt systematiska, elementfokuserade angreppss\"att: varje
    kapitel \"ar organiserat kring fr\aa{}gan ``vad g\"or vanadin
    specifikt, och hur skiljer sig detta fr\aa{}n Nb och Ti?''
  \item Sin integration av grundl\"aggande termodynamik och
    utskiljningsteori med fullskaliga industridata fr\aa{}n svenska,
    koreanska, finska, italienska, kinesiska och brasilianska st\aa{}lverk.
  \item Sin f\"orespr\aa{}kan f\"or kv\"ave som ett medvetet tillagt
    legeringselement---en st\aa{}ndpunkt som var kontroversiell n\"ar den
    f\"orst framf\"ordes av Swerea-gruppen p\aa{} 1990-talet men som
    sedan dess har validerats brett industriellt.
  \item Sin utveckling av kvantitativa process--egenskapsmodeller (MICDEL)
    som \"overbryggar klyftan mellan laboratorief\"orst\aa{}else och
    industriell praxis.
\end{itemize}

Monografin representerar kulmen av approximately 30 \aa{}rs
samarbete inom SIMR/Swerea KIMAB:s mikrolegeringsgrupp, d\"ar
Hutchinson var en central medlem under hela perioden.

%% ===================================================================
\section*{Sammanfattning}

Hutchinsons bidrag kan grupperas i flera sammankopplade teman:

\begin{enumerate}
  \item \textbf{Textur--egenskapssamband:} att uppst\"alla kvantitativa
    samband mellan kristallografisk textur (s\"arskilt
    $\gamma$-fibern i BCC och kubkomponenten i FCC) och makroskopisk
    mekanisk anisotropi ($\bar{r}$-v\"arden, \"oronbildning).

  \item \textbf{Rekristallisationsmekanismer:} att identifiera de
    deformationssubstrukturer (skjuvband, \"overg\aa{}ngsband,
    kubband) som fungerar som nukleeringsst\"allen, och klar\-l\"agga
    rollerna f\"or orienterad nukleering och selektiv tillv\"axt.

  \item \textbf{Process--texturstyrning:} att \"overs\"atta mekanistisk
    f\"orst\aa{}else till praktisk v\"agledning f\"or industriell
    bearbetning---varmvalsning, kallvalsning och
    gl\"odgningsscheman---f\"or att uppn\aa{} m\aa{}ltexter i st\aa{}l
    och aluminiumlegeringar.

  \item \textbf{Mikrolegerade st\aa{}l:} att klar\-l\"agga de distinkta
    rollerna f\"or V, Nb och Ti i termomekanisk bearbetning, med
    s\"arskilda bidrag till f\"orst\aa{}elsen av vanadinutskiljning och
    rekristallisationskontrollerad valsning.

  \item \textbf{Elst\aa{}l:} att bidra till f\"orst\aa{}elsen av
    Goss-texturselektion vid sekund\"ar rekristallisation i
    kornorienterade Fe--Si-st\aa{}l.

  \item \textbf{Process\"overvakning:} att vara pioneer inom
    till\"ampning av laserultraljud f\"or realtidsm\"atning av
    mikrostruktur under varmvalsning av bandpl\aa{}t.
\end{enumerate}

%% ===================================================================
\begin{thebibliography}{99}

\bibitem{Hutchinson1984IMR}
W.B.~Hutchinson,
``Development and control of annealing textures in low-carbon steels,''
\emph{International Materials Reviews}, vol.~29, no.~1, pp.~25--72, 1984.

\bibitem{Hutchinson2007MSF}
W.B.~Hutchinson,
``Deformation substructures and recrystallisation,''
\emph{Materials Science Forum}, vols.~558--559, pp.~13--22, 2007.

\bibitem{Ridha1982}
A.A.~Ridha and W.B.~Hutchinson,
``Recrystallisation mechanisms and the origin of cube texture in copper,''
\emph{Acta Metallurgica}, vol.~30, pp.~1929--1939, 1982.

\bibitem{Homma2003}
H.~Homma and B.~Hutchinson,
``Orientation dependence of secondary recrystallisation in silicon--iron,''
\emph{Acta Materialia}, vol.~51, pp.~3795--3805, 2003.

\bibitem{Lagneborg2014}
R.~Lagneborg, T.~Siwecki, S.~Zajac, and B.~Hutchinson,
\emph{The Role of Vanadium in Microalloyed Steels},
Swerea KIMAB Report KIMAB-2014-115, June 2014.
(Originally published in \emph{Scandinavian Journal of Metallurgy},
vol.~28, pp.~186--241, 1999; expanded edition available from Vanitec.)

\bibitem{Oscarsson1991}
A.~Oscarsson, W.B.~Hutchinson, and H.-E.~Ekstr\"om,
``Influence of initial microstructure on texture and earing in aluminium
sheet after cold rolling and annealing,''
\emph{Materials Science and Technology}, vol.~7, pp.~554--564, 1991.

\bibitem{Hutchinson2015}
B.~Hutchinson,
``Critical Assessment~16: Anisotropy in metals,''
\emph{Materials Science and Technology}, vol.~31, no.~12,
pp.~1393--1401, 2015.

\bibitem{Zajac1991}
S.~Zajac, T.~Siwecki, B.~Hutchinson, and M.~Attleg\aa{}rd,
``Recrystallization controlled rolling and accelerated cooling for high
strength and toughness in V--Ti--N steels,''
\emph{Metallurgical Transactions~A}, vol.~22A, pp.~2681--2694, 1991.

\bibitem{Hutchinson2011LUS}
B.~Hutchinson et~al.,
``Laser ultrasonics for process control in the metal industry,''
\emph{Nondestructive Testing and Evaluation}, vol.~26, nos.~3--4, 2011.

\bibitem{Hutchinson1999Phil}
B.~Hutchinson,
``Deformation microstructures and textures in steels,''
\emph{Philosophical Transactions of the Royal Society~A}, vol.~357,
no.~1756, pp.~1471--1485, 1999.

\bibitem{Hutchinson2025pssb}
B.~Hutchinson, M.~Malmstr\"om, A.~Jansson, and J.~L\"onnqvist,
``An experimental study of elasticity in polycrystalline iron,''
\emph{physica status solidi~(b)}, vol.~262, no.~6, 2025.

\bibitem{Hatherly1979}
M.~Hatherly and W.B.~Hutchinson,
\emph{An Introduction to Textures in Metals},
Institute of Metals, London, 1979; 2nd~ed.\ 1990.

\bibitem{Hutchinson2006Met}
B.~Hutchinson and J.~Hagstr\"om,
``Austenite decomposition structures in the Gibeon meteorite,''
\emph{Metallurgical and Materials Transactions~A}, vol.~37,
pp.~1811--1818, 2006.

\bibitem{Ray2011}
R.K.~Ray, B.~Hutchinson, and C.~Ghosh,
``Back-annealing of cold rolled steels through recovery and/or partial
recrystallisation,''
\emph{International Materials Reviews}, vol.~56, no.~2, pp.~73--97, 2011.

\bibitem{LagneborgZajac2001}
R.~Lagneborg and S.~Zajac,
``A model for interphase precipitation in V-microalloyed structural steels,''
\emph{Metallurgical and Materials Transactions~A}, vol.~32A, pp.~39--46,
2001.

\bibitem{Siwecki1995}
T.~Siwecki, B.~Hutchinson, and S.~Zajac,
``Recrystallisation controlled rolling of HSLA steels,''
Proc.\ Conf.\ Microalloying~'95, Pittsburgh, Iron \& Steel Soc., 1995,
pp.~197--211.

\bibitem{Siwecki1992}
T.~Siwecki and B.~Hutchinson,
``Modelling of microstructure evolution during recrystallisation
controlled rolling of HSLA steels,''
Proc.\ 33rd Mechanical Working and Steel Processing Conf., Warrendale,
ISS-AIME, 1992, pp.~397--406.

\bibitem{Hutchinson2014Bainite}
B.~Hutchinson, J.~Hagstr\"om, T.~Siwecki, J.~Komenda, R.~Lagneborg,
J.-E.~Hedin, and M.~Gladh,
``New vanadium-microalloyed bainitic 700~MPa strip steel product,''
\emph{Ironmaking and Steelmaking}, vol.~41, pp.~1--6, 2014.

\end{thebibliography}

\end{document}
